\GAGASection{\S 3. Correspondence between algebraic sheaves and coherent analytic sheaves}{gaga:section:3}

\GAGASubsection{9. The analytic sheaf associated with an algebraic sheaf}{gaga:no:9}

Let $\mathrm{X}$ be an algebraic variety, and let $\mathrm{X}^h$ be the analytic space associated with it by the construction of \hyperref[gaga:no:5]{no.~5}. If $\mathcal{F}$ is any sheaf on $\mathrm{X}$, we shall endow the set $\mathcal{F}$ with a new topology that makes it a sheaf on $\mathrm{X}^h$. This topology is defined as follows. If $\pi:\mathcal{F}\to\mathrm{X}$ denotes the projection of $\mathcal{F}$ onto $\mathrm{X}$, embed $\mathcal{F}$ in $\mathrm{X}^h\times\mathcal{F}$ by the map $f\mapsto(\pi(f),f)$, and give $\mathcal{F}$ the topology induced by that of $\mathrm{X}^h\times\mathcal{F}$. It is immediately verified that the set $\mathcal{F}$ has thereby been endowed with the structure of a sheaf on $\mathrm{X}^h$; we shall denote this sheaf by $\mathcal{F}'$. Thus, for every $x\in\mathrm{X}$, we have $\mathcal{F}'_x=\mathcal{F}_x$; the sheaves $\mathcal{F}$ and $\mathcal{F}'$ differ only in their topology ($\mathcal{F}'$ is simply the \emph{inverse image} of $\mathcal{F}$ under the continuous map $\mathrm{X}^h\to\mathrm{X}$).

The foregoing applies in particular to the sheaf $\mathcal{O}$ of local rings of $\mathrm{X}$. Proposition~\hyperref[gaga:prop:3]{3} of \hyperref[gaga:no:6]{no.~6} allows us to identify the sheaf $\mathcal{O}'$ thereby obtained with a subsheaf of the sheaf $\mathcal{H}$ of germs of holomorphic functions on $\mathrm{X}^h$.

\GAGATarget{gaga:def:2}\textsc{Definition 2}. --- \emph{Let $\mathcal{F}$ be an algebraic sheaf on $\mathrm{X}$. The analytic sheaf associated with $\mathcal{F}$ is the sheaf $\mathcal{F}^h$ on $\mathrm{X}^h$ defined by}

\begin{equation*}
\mathcal{F}^h=\mathcal{F}'\otimes\mathcal{H},
\end{equation*}

\emph{where the tensor product is taken over the sheaf of rings $\mathcal{O}'$.}

(In other words, $\mathcal{F}^h$ is obtained from $\mathcal{F}'$ by extending the ring of operators to $\mathcal{H}$.)

The sheaf $\mathcal{F}^h$ is a sheaf of $\mathcal{H}$-modules, that is, an \emph{analytic} sheaf. The injection $\mathcal{O}'\to\mathcal{H}$ defines a canonical homomorphism $\alpha:\mathcal{F}'\to\mathcal{F}^h$.

Every algebraic homomorphism (that is, every $\mathcal{O}$-linear homomorphism)

\begin{equation*}
\varphi:\mathcal{F}\longrightarrow\mathcal{G}
\end{equation*}

defines, by extension of the ring of operators, an analytic homomorphism

\begin{equation*}
\varphi^h:\mathcal{F}^h\longrightarrow\mathcal{G}^h.
\end{equation*}

Thus $\mathcal{F}^h$ is a \emph{covariant functor} of $\mathcal{F}$.

\GAGATarget{gaga:prop:10}\textsc{Proposition 10}. --- \GAGATarget{gaga:prop:10:item:a}a) \emph{The functor $\mathcal{F}^h$ is exact.}

\GAGATarget{gaga:prop:10:item:b}b) \emph{For every algebraic sheaf $\mathcal{F}$, the homomorphism $\alpha:\mathcal{F}'\to\mathcal{F}^h$ is injective.}

\GAGATarget{gaga:prop:10:item:c}c) \emph{If $\mathcal{F}$ is a coherent algebraic sheaf, then $\mathcal{F}^h$ is a coherent analytic sheaf.}

If $\mathcal{F}_1\to\mathcal{F}_2\to\mathcal{F}_3$ is an exact sequence of algebraic sheaves, then the same is plainly true of $\mathcal{F}'_1\to\mathcal{F}'_2\to\mathcal{F}'_3$, and hence also of the sequence

\begin{equation*}
\mathcal{F}'_1\otimes\mathcal{H}\longrightarrow
\mathcal{F}'_2\otimes\mathcal{H}\longrightarrow
\mathcal{F}'_3\otimes\mathcal{H},
\end{equation*}

by Corollary~\hyperref[gaga:cor:1]{1} to Proposition~\hyperref[gaga:prop:3]{3}; this proves a). Assertion b) also follows from the same corollary.

To prove c), first observe that $\mathcal{O}^h=\mathcal{H}$. If $\mathcal{F}$ is algebraically coherent and $x$ is a point of $X$, one can find an exact sequence

\begin{equation*}
\mathcal{O}^q\longrightarrow\mathcal{O}^p\longrightarrow\mathcal{F}\longrightarrow0
\end{equation*}

valid in a Z-neighbourhood $U$ of $x$. By a), this yields an exact sequence

\begin{equation*}
\mathcal{H}^q\longrightarrow\mathcal{H}^p\longrightarrow\mathcal{F}^h\longrightarrow0
\end{equation*}

valid on $U$. Since $U$ is a neighbourhood of $x$ and the sheaf of rings $\mathcal{H}$ is coherent (Proposition~\hyperref[gaga:prop:1]{1}, \hyperref[gaga:no:3]{no.~3}), this proves that $\mathcal{F}^h$ is coherent (\GAGACite{18}, no.~15).

The preceding proposition shows in particular that, if $\mathcal{J}$ is a sheaf of ideals of $\mathcal{O}$, then $\mathcal{J}^h$ is precisely the sheaf of ideals of $\mathcal{H}$ generated by the elements of $\mathcal{J}$.

\GAGASubsection{10. Extension of a sheaf}{gaga:no:10}

Let $Y$ be a Z-closed subvariety of the algebraic variety $X$, and let $\mathcal{F}$ be a coherent algebraic sheaf on $Y$. Denote by $\mathcal{F}^X$ the sheaf obtained by extending $\mathcal{F}$ by zero on $X-Y$ (cf. \GAGACite{18}, no.~5). We know that $\mathcal{F}^X$ is a coherent algebraic sheaf on $X$, and the sheaf $(\mathcal{F}^X)^h$ is therefore defined; it is a coherent analytic sheaf on $X^h$. On the other hand, $\mathcal{F}^h$ is a coherent analytic sheaf on $Y^h$, and it may be extended by zero on $X^h-Y^h$, thereby giving a new sheaf $(\mathcal{F}^h)^X$. We have:

\GAGATarget{gaga:prop:11}\textsc{Proposition 11}. --- \emph{The sheaves $(\mathcal{F}^h)^X$ and $(\mathcal{F}^X)^h$ are canonically isomorphic.}

Both sheaves in question vanish outside $Y^h$; it will therefore suffice to show that their restrictions to $Y^h$ are isomorphic.

Let $x$ be a point of $Y$. To simplify the notation, put

\begin{equation*}
A=\mathcal{O}_{x,\,X},\qquad A'=\mathcal{O}_{x,\,Y},\qquad
B=\mathcal{H}_{x,\,X},\qquad B'=\mathcal{H}_{x,\,Y},\qquad E=\mathcal{F}_x.
\end{equation*}

Then

\begin{equation*}
(\mathcal{F}^h)^X_x=E\otimes_{A'}B'
\qquad\text{and}\qquad
(\mathcal{F}^X)^h_x=E\otimes_AB.
\end{equation*}

The ring $A'$ is the quotient of $A$ by an ideal $\mathfrak{a}$, and Proposition~\hyperref[gaga:prop:4]{4} of \hyperref[gaga:no:6]{no.~6} gives $B'=B/\mathfrak{a}B=B\otimes_AA'$. By associativity of the tensor product, we therefore obtain an isomorphism

\begin{equation*}
\theta_x:E\otimes_{A'}B'=E\otimes_{A'}A'\otimes_AB\longrightarrow E\otimes_AB,
\end{equation*}

which varies continuously with $x$, as is readily seen; the proposition follows.

Proposition~\hyperref[gaga:prop:11]{11} may be expressed by saying that the functor $\mathcal{F}^h$ is \emph{compatible with the usual identification of $\mathcal{F}$ with $\mathcal{F}^X$}.

\GAGASubsection{11. Homomorphisms induced on cohomology}{gaga:no:11}

Use the notation of \hyperref[gaga:no:9]{no.~9}. Let $X$ be an algebraic variety, $\mathcal{F}$ an algebraic sheaf on $X$, and $\mathcal{F}^h$ the analytic sheaf associated with $\mathcal{F}$. If $U$ is a Z-open subset of $X$ and $s$ is a section of $\mathcal{F}$ over $U$, we may regard $s$ as a section $s'$ of $\mathcal{F}'$ over the open subset $U^h$ of $X^h$; then $\alpha(s')=s'\otimes1$ is a section of $\mathcal{F}^h=\mathcal{F}'\otimes\mathcal{H}$ over $U^h$. The map $s\mapsto\alpha(s')$ is a homomorphism

\begin{equation*}
\varepsilon:\Gamma(U,\mathcal{F})\longrightarrow\Gamma(U^h,\mathcal{F}^h).
\end{equation*}

Now let $\mathfrak{U}=\{U_i\}$ be a finite Z-open covering of $X$. The $U_i^h$ form a finite open covering of $X^h$, which we shall denote by $\mathfrak{U}^h$. For every system of indices $i_0,\ldots,i_q$, the preceding construction gives a canonical homomorphism

\begin{equation*}
\varepsilon:\Gamma(U_{i_0}\cap\cdots\cap U_{i_q},\mathcal{F})
\longrightarrow
\Gamma(U^h_{i_0}\cap\cdots\cap U^h_{i_q},\mathcal{F}^h),
\end{equation*}

and hence a homomorphism

\begin{equation*}
\varepsilon:C(\mathfrak{U},\mathcal{F})\longrightarrow C(\mathfrak{U}^h,\mathcal{F}^h),
\end{equation*}

with the notation of \GAGACite{18}, no.~18.

This homomorphism commutes with the coboundary $d$ and therefore induces, on passing to cohomology, homomorphisms

\begin{equation*}
\varepsilon:\mathrm{H}^q(\mathfrak{U},\mathcal{F})
\longrightarrow
\mathrm{H}^q(\mathfrak{U}^h,\mathcal{F}^h).
\end{equation*}

Finally, on passing to the inductive limit over $\mathfrak{U}$, we obtain \emph{the homomorphisms induced on the cohomology groups}

\begin{equation*}
\varepsilon:\mathrm{H}^q(X,\mathcal{F})
\longrightarrow
\mathrm{H}^q(X^h,\mathcal{F}^h).
\end{equation*}

These homomorphisms have the usual functorial properties; they commute with homomorphisms $\varphi:\mathcal{F}\to\mathcal{G}$. If

\begin{equation*}
0\longrightarrow\mathcal{A}\longrightarrow\mathcal{B}\longrightarrow\mathcal{C}\longrightarrow0
\end{equation*}

is an exact sequence of algebraic sheaves in which $\mathcal{A}$ is \emph{coherent}, the diagram

\begin{center}
\begin{tikzcd}
\mathrm{H}^q(X,\mathcal{C}) \arrow[r,"\delta"] \arrow[d,"\varepsilon"'] &
\mathrm{H}^{q+1}(X,\mathcal{A}) \arrow[d,"\varepsilon"'] \\
\mathrm{H}^q(X^h,\mathcal{C}^h) \arrow[r,"\delta"] &
\mathrm{H}^{q+1}(X^h,\mathcal{A}^h)
\end{tikzcd}
\end{center}

is commutative. This may be seen, for example, by taking the coverings $\mathfrak{U}$ to be coverings by affine open sets (cf. \GAGACite{18}).

\GAGASubsection{12. Projective varieties. Statement of the theorems}{gaga:no:12}

Suppose that $\mathrm{X}$ is a \emph{projective variety}, that is, a Z-closed subvariety of a projective space $\mathbf{P}_r(\mathbf{C})$. We then have the following theorems, which will be proved in the remainder of this section.

\GAGATarget{gaga:thm:1}\textsc{Theorem 1}. --- \emph{For every coherent algebraic sheaf $\mathcal{F}$ on $\mathrm{X}$ and every integer $q\geqslant0$, the homomorphism}

\begin{equation*}
\varepsilon:\mathrm{H}^q(\mathrm{X},\mathcal{F})
\longrightarrow
\mathrm{H}^q(\mathrm{X}^h,\mathcal{F}^h),
\end{equation*}

\emph{defined in \hyperref[gaga:no:11]{no.~11}, is bijective.}

For $q=0$, in particular, this gives an isomorphism from $\Gamma(\mathrm{X},\mathcal{F})$ onto $\Gamma(\mathrm{X}^h,\mathcal{F}^h)$.

\GAGATarget{gaga:thm:2}\textsc{Theorem 2}. --- \emph{If $\mathcal{F}$ and $\mathcal{G}$ are two coherent algebraic sheaves on $\mathrm{X}$, every analytic homomorphism from $\mathcal{F}^h$ to $\mathcal{G}^h$ arises from one and only one algebraic homomorphism from $\mathcal{F}$ to $\mathcal{G}$.}

\GAGATarget{gaga:thm:3}\textsc{Theorem 3}. --- \emph{For every coherent analytic sheaf $\mathcal{M}$ on $\mathrm{X}^h$, there exists a coherent algebraic sheaf $\mathcal{F}$ on $\mathrm{X}$ such that $\mathcal{F}^h$ is isomorphic to $\mathcal{M}$. Moreover, this property determines $\mathcal{F}$ uniquely up to isomorphism.}

\emph{Remarks}. --- 1. These three theorems mean that the theory of coherent analytic sheaves on $\mathrm{X}^h$ essentially coincides with that of coherent algebraic sheaves on $\mathrm{X}$. They of course depend on $\mathrm{X}$ being \emph{projective}, and fail even for an affine variety.

2. The homomorphism $\varepsilon$ may be factored as

\begin{equation*}
\mathrm{H}^q(\mathrm{X},\mathcal{F})
\longrightarrow
\mathrm{H}^q(\mathrm{X}^h,\mathcal{F}')
\longrightarrow
\mathrm{H}^q(\mathrm{X}^h,\mathcal{F}^h).
\end{equation*}

Theorem~\hyperref[gaga:thm:1]{1} raises the question whether $\mathrm{H}^q(\mathrm{X},\mathcal{F})\to\mathrm{H}^q(\mathrm{X}^h,\mathcal{F}')$ is bijective. The answer is negative. Indeed, if this homomorphism were bijective for every coherent algebraic sheaf $\mathcal{F}$, it would also be so for the constant sheaf $\mathrm{K}=\mathbf{C}(\mathrm{X})$ of rational functions on $\mathrm{X}$ (assumed irreducible), since that sheaf is the union of coherent sheaves (compare \GAGACite{19}, \S 2). But $\mathrm{H}^q(\mathrm{X},\mathrm{K})=0$ for $q>0$, whereas $\mathrm{H}^q(\mathrm{X}^h,\mathrm{K})$ is a $\mathrm{K}$-vector space whose dimension equals the $q$-th Betti number of $\mathrm{X}^h$.
